% =============================================================================
%  Residual Reinforcement Learning for SEA-Actuated Manipulator
%  Implementation notes: Gymnasium env, PPO training, evaluation pipeline
% =============================================================================

\documentclass[11pt, a4paper]{article}

\usepackage{amsmath, amssymb, bm}
\usepackage{geometry}
\usepackage{hyperref}
\usepackage{booktabs}
\usepackage{xcolor}
\usepackage{listings}
\usepackage{cleveref}
\usepackage{tikz}
\usepackage{enumitem}
\usepackage[backgroundcolor=blue!5, linecolor=blue!40, linewidth=1pt,
            innertopmargin=6pt, innerbottommargin=6pt,
            innerleftmargin=7pt, innerrightmargin=7pt]{mdframed}

\geometry{margin=2.2cm}

\lstset{
  language=Python,
  basicstyle=\ttfamily\small,
  keywordstyle=\color{blue!70!black},
  commentstyle=\color{gray},
  stringstyle=\color{orange!80!black},
  frame=single,
  breaklines=true,
  columns=fullflexible,
  backgroundcolor=\color{gray!8},
  xleftmargin=6pt, xrightmargin=6pt,
  aboveskip=4pt, belowskip=4pt,
}

% ── Standard project macros ──────────────────────────────────────────────────
\newcommand{\bq}{\bm{q}}
\newcommand{\bp}{\bm{p}}
\newcommand{\bu}{\bm{u}}
\newcommand{\btau}{\bm{\tau}}
\newcommand{\bF}{\bm{F}}
\newcommand{\bJ}{\bm{J}}
\newcommand{\bM}{\bm{M}}
\newcommand{\bC}{\bm{C}}
\newcommand{\bG}{\bm{G}}

\title{\textbf{Residual RL for SEA-Actuated Manipulator}\\
       \large \texttt{rl/train\_ppo\_residual.py} \;|\;
              \texttt{rl/envs/manipulator\_residual\_env.py} \;|\;
              \texttt{rl/eval\_ppo\_residual.py}}
\date{April 2026}
\author{Isaac-sim robotics notes}

\begin{document}
\maketitle
\tableofcontents
\newpage

% =============================================================================
\section{Motivation}
% =============================================================================

The 2-DOF cable-driven manipulator uses a computed-torque (CT) controller
to track end-effector trajectories.  Joint~1 is rigidly actuated while
joint~2 is driven through a Series Elastic Actuator (SEA) cable, introducing:

\begin{itemize}[nosep]
  \item \textbf{Phase lag} — the motor servo has finite bandwidth $\omega_m$,
        so the spring-mediated torque $\tau_2^{\text{sea}}$ lags behind the
        CT-desired torque $\tau_2^{\text{ct}}$.
  \item \textbf{Model mismatch} — the CT law relies on $\bM(\bq)$ and
        $\bm{h}(\bq,\dot{\bq})$ from the ArticulationView, which may
        differ from the true PhysX dynamics under fast motions.
  \item \textbf{Spring compliance} — the cable spring ($k_s$) absorbs
        energy, causing steady-state tracking offsets that grow with
        trajectory curvature.
\end{itemize}

\begin{mdframed}
\textbf{Key idea (Residual RL).}\;
Rather than replacing the CT controller entirely, we learn a \emph{small}
additive correction $\Delta\btau$ that is summed with the CT output
\emph{before} the SEA model.  This preserves the well-understood baseline
while letting the RL agent compensate for the phenomena above.
\end{mdframed}


% =============================================================================
\section{System Architecture}
% =============================================================================

The full control loop is:

\begin{center}
\begin{tikzpicture}[
  block/.style={draw, minimum width=2.0cm, minimum height=0.9cm,
                rounded corners=2pt, font=\small},
  arr/.style={->, thick},
  node distance=1.6cm and 1.8cm,
]
  \node[block] (ct)   {CT Controller};
  \node[block, below=0.8cm of ct] (rl) {RL Policy $\pi_\theta$};
  \node[block, right=1.6cm of ct] (add) {$+$};
  \node[block, right=1.4cm of add] (sea) {SEA Cable};
  \node[block, right=1.4cm of sea] (plant) {PhysX Plant};

  \draw[arr] (ct.east) -- node[above,font=\footnotesize]{$\btau^{\text{ct}}$} (add.west);
  \draw[arr] (rl.east) -| node[near start,below,font=\footnotesize]{$\Delta\btau$} (add.south);
  \draw[arr] (add.east) -- node[above,font=\footnotesize]{$\btau^{\text{ct}}\!+\!\Delta\btau$} (sea.west);
  \draw[arr] (sea.east) -- node[above,font=\footnotesize]{$\btau^{\text{sea}}$} (plant.west);

  % Feedback
  \draw[arr] (plant.south) -- ++(0,-0.8) -| node[near start,below,font=\footnotesize]{observation $\bm{o}_t$} (rl.south);
  \draw[arr] (plant.south) -- ++(0,-0.8) -| (ct.south);
\end{tikzpicture}
\end{center}

At every timestep $\Delta t = 0.01\;\text{s}$ (100\,Hz):
\begin{enumerate}[nosep]
  \item Read joint state $(\bq, \dot{\bq})$ from the PhysX plant.
  \item Trajectory source provides $(\bp_{\text{ref}}, \dot{\bp}_{\text{ref}}, \ddot{\bp}_{\text{ref}})$ in task space.
  \item IK + analytical Jacobian convert to joint-space references $(\bq_{\text{des}}, \dot{\bq}_{\text{ref}}, \ddot{\bq}_{\text{ref}})$.
  \item CT computes $\btau^{\text{ct}} = \bM(\bq)\,\bm{a}_{\text{des}} + \bm{h}(\bq,\dot{\bq})$.
  \item RL policy outputs $\Delta\btau = \pi_\theta(\bm{o}_t)$.
  \item SEA model processes $\btau^{\text{ct}} + \Delta\btau$ through the motor servo + spring-damper.
  \item Resulting $\btau^{\text{sea}}$ is applied to the PhysX articulation.
\end{enumerate}


% =============================================================================
\section{Gymnasium Environment}
\label{sec:env}
% =============================================================================

The environment is implemented in
\texttt{rl/envs/manipulator\_residual\_env.py} as a
\texttt{gymnasium.Env} subclass.

\subsection{Design Decision: External Scene Creation}

\begin{mdframed}
The environment does \textbf{not} create its own \texttt{SimulationApp} or
\texttt{World}.  Instead it receives pre-built Isaac Sim objects (robot,
world, controller, SEA) and drives the existing simulation.  This avoids
the overhead of re-creating the USD stage on every \texttt{env.reset()} and
allows the same scene to be shared between training and evaluation scripts.
\end{mdframed}

Constructor signature:
\begin{lstlisting}
ManipulatorResidualEnv(
    robot,            # CupManipulatorTendonIsaac (initialized)
    world,            # omni.isaac.core.World
    ct_controller,    # ComputedTorqueController
    sea_actuator,     # SEACableActuatorNP
    traj_source,      # PreambleTrajectorySource
    L1, L2,           # link lengths [m]
    dt=0.01,          # physics timestep
    max_episode_steps=2000,
    residual_max=5.0, # action bounds [Nm]
    reward_weights={"tracking": 100, "effort": 0.01, "smoothness": 0.001},
    solve_ik_fn, fk_fn, ik_to_jt_fn,  # kinematics functions
)
\end{lstlisting}

\subsection{Observation Space (10-D)}

\begin{center}
\renewcommand{\arraystretch}{1.15}
\begin{tabular}{clcc}
\toprule
\textbf{Index} & \textbf{Quantity} & \textbf{Symbol} & \textbf{Bounds} \\
\midrule
0 & Joint 1 position  & $q_1$         & $[-\pi, \pi]$ \\
1 & Joint 2 position  & $q_2$         & $[-\pi, \pi]$ \\
2 & Joint 1 velocity  & $\dot{q}_1$   & $[-10, 10]$ \\
3 & Joint 2 velocity  & $\dot{q}_2$   & $[-10, 10]$ \\
4 & EE error X        & $e_x$         & $[-0.5, 0.5]$ \\
5 & EE error Y        & $e_y$         & $[-0.5, 0.5]$ \\
6 & Spring extension  & $\delta$      & $[-0.1, 0.1]$ \\
7 & Cable force       & $F_{\text{cable}}$ & $[-500, 500]$ \\
8 & CT desired $\tau_2$ & $\tau_2^{\text{ct}}$ & $[-50, 50]$ \\
9 & SEA applied $\tau_2$ & $\tau_2^{\text{sea}}$ & $[-50, 50]$ \\
\bottomrule
\end{tabular}
\end{center}

\textbf{Design rationale:}
\begin{itemize}[nosep]
  \item $e_x, e_y$ tell the agent \emph{how far off} the EE is (task-space error).
  \item $\delta$ and $F_{\text{cable}}$ expose the SEA internal state so the
        agent can anticipate spring dynamics.
  \item $\tau_2^{\text{ct}}$ vs $\tau_2^{\text{sea}}$ reveal the CT--SEA
        mismatch the agent should compensate.
\end{itemize}

\subsection{Action Space (2-D)}

\[
  \bm{a}_t = \Delta\btau = [\Delta\tau_1,\; \Delta\tau_2]
  \;\in\; [-\delta_{\max},\; \delta_{\max}]^2
\]

Default $\delta_{\max} = 5\;\text{Nm}$.  This is intentionally small
relative to the CT output ($|\tau^{\text{ct}}| \leq 50\;\text{Nm}$) so
the RL agent acts as a \emph{correction}, not a replacement.

\subsection{Reward Function}

\begin{align}
  r_t &= \underbrace{-\alpha \|\bm{e}_t\|^2}_{%
    \text{tracking}} \;
    \underbrace{-\;\beta \|\Delta\btau_t\|^2}_{%
    \text{effort}} \;
    \underbrace{-\;\gamma \|\btau_t - \btau_{t-1}\|^2}_{%
    \text{smoothness}}
  \label{eq:reward}
\end{align}

where $\bm{e}_t = \bp_{\text{ee}} - \bp_{\text{ref}}$ is the task-space
tracking error.  Default weights:

\begin{center}
\begin{tabular}{lrl}
\toprule
\textbf{Component} & \textbf{Weight} & \textbf{Purpose} \\
\midrule
$\alpha$ (tracking)   & 100.0   & Penalize EE deviation \\
$\beta$ (effort)      & 0.01    & Regularize residual magnitude \\
$\gamma$ (smoothness) & 0.001   & Penalize torque jerk \\
\bottomrule
\end{tabular}
\end{center}

$\beta$ and $\gamma$ prevent the policy from learning large, noisy
corrections that would stress the cable or excite vibrations.

\subsection{Episode Reset}

On \texttt{env.reset()}:
\begin{enumerate}[nosep]
  \item Joint angles are set to nominal $\pm$ uniform noise $\mathcal{U}(-0.05, 0.05)$\;rad
        (domain randomization of initial conditions).
  \item Velocities zeroed.
  \item SEA motor cable position re-initialized: $l_m = r_p \cdot q_2$.
  \item IK seed computed at first trajectory waypoint.
  \item One physics step to settle.
\end{enumerate}

\subsection{Termination Conditions}

\begin{itemize}[nosep]
  \item \textbf{Terminated}: $|q_i| > 150°$ (joint limit violation).
  \item \textbf{Truncated}: step count $\geq$ \texttt{max\_episode\_steps}
        (default 2000 = 20\,s).
\end{itemize}


% =============================================================================
\section{PPO Training}
\label{sec:train}
% =============================================================================

Training is handled by \texttt{rl/train\_ppo\_residual.py} using
Stable-Baselines3 (SB3) v2.8.

\subsection{Isaac Sim Scene Construction}

The training script handles the full Isaac Sim lifecycle:

\begin{enumerate}[nosep]
  \item Pre-parse \texttt{--render} from \texttt{sys.argv} (before argparse).
  \item Create \texttt{SimulationApp} with headless mode.
  \item Build \texttt{World}, load URDF $\to$ USD, weld base, set properties.
  \item Initialize \texttt{ArticulationView} for dynamics queries ($\bM$, $\bm{h}$).
  \item Create trajectory, CT controller, and SEA actuator.
  \item Wrap everything in \texttt{ManipulatorResidualEnv} $\to$ \texttt{Monitor}.
\end{enumerate}

\begin{mdframed}
\textbf{Important:} \texttt{SimulationApp} must be instantiated \emph{before}
any \texttt{omni.*} imports.  The \texttt{--render} flag is pre-parsed with
a simple \texttt{sys.argv} loop because \texttt{argparse} cannot run before
the app is created.
\end{mdframed}

\subsection{PPO Hyperparameters}

\begin{center}
\renewcommand{\arraystretch}{1.15}
\begin{tabular}{lrl}
\toprule
\textbf{Parameter} & \textbf{Default} & \textbf{CLI flag} \\
\midrule
Learning rate       & $3 \times 10^{-4}$ & \texttt{--learning-rate} \\
Rollout steps       & 2048               & \texttt{--n-steps} \\
Mini-batch size     & 64                 & \texttt{--batch-size} \\
Epochs per update   & 10                 & \texttt{--n-epochs} \\
Discount $\gamma$   & 0.99               & \texttt{--gamma} \\
Clip range          & 0.2                & \texttt{--clip-range} \\
Entropy coeff.      & 0.0                & \texttt{--ent-coef} \\
Network arch.       & $[128,\,128]$      & (hardcoded) \\
Device              & CPU                & (hardcoded) \\
\bottomrule
\end{tabular}
\end{center}

\textbf{Why CPU?}\; Isaac Sim owns the GPU for PhysX rendering.
Running the policy on CPU avoids CUDA context conflicts.  The 2-layer MLP
with 128 neurons is small enough that CPU inference is not a bottleneck.

\subsection{Network Architecture}

Both the policy $\pi_\theta$ and value function $V_\phi$ use separate
128--128 MLPs:

\begin{lstlisting}
policy_kwargs = dict(
    net_arch=dict(pi=[128, 128], vf=[128, 128]),
)
\end{lstlisting}

With a 10-D input and 2-D output, the policy has approximately:
\[
  (10 \times 128 + 128) + (128 \times 128 + 128) + (128 \times 2 + 2)
  = 18\,178 \;\text{parameters}
\]

\subsection{Checkpointing and Logging}

\begin{itemize}[nosep]
  \item \textbf{Checkpoints}: saved every 10\,000 steps to \texttt{rl/checkpoints/}
        (configurable via \texttt{--save-freq}).
  \item \textbf{Logging}: stdout + CSV + TensorBoard to \texttt{rl/logs/}.
  \item \textbf{Resume}: \texttt{--resume rl/checkpoints/ppo\_residual\_latest}
        reloads weights and continues training.
  \item \textbf{SIGINT}: Ctrl+C triggers a clean save before exit
        (using \texttt{signal.signal}, not \texttt{try/except}).
\end{itemize}

\subsection{Training Flow}

\begin{enumerate}[nosep]
  \item SB3 collects a rollout of $n_{\text{steps}} = 2048$ transitions by
        calling \texttt{env.step(action)} repeatedly.
  \item Rollout buffer is split into mini-batches of size 64.
  \item PPO objective is optimized for 10 epochs with clipped surrogate loss.
  \item Repeat until \texttt{total\_timesteps} reached or Ctrl+C.
\end{enumerate}

Each rollout collects $\approx 20$\,s of wall time at headless RTF $\approx 2$--$5\times$.


% =============================================================================
\section{Evaluation}
\label{sec:eval}
% =============================================================================

\texttt{rl/eval\_ppo\_residual.py} loads a trained checkpoint and runs
deterministic rollouts for analysis.

\subsection{Evaluation Modes}

\begin{center}
\begin{tabular}{lp{8cm}}
\toprule
\textbf{Flag} & \textbf{Behavior} \\
\midrule
(default) & Load model, run one episode with CT+RL, plot results. \\
\texttt{--compare-baseline} & Additionally run an episode with $\Delta\btau = \bm{0}$
    (CT-only) for side-by-side comparison. \\
\texttt{--render native} & Visualize in real-time Isaac Sim window. \\
\bottomrule
\end{tabular}
\end{center}

\subsection{Output Plots}

The evaluation script produces a multi-panel figure saved to
\texttt{plots/eval\_ppo\_residual\_<timestamp>.png}:

\begin{enumerate}[nosep]
  \item \textbf{Tracking error} [mm] vs time --- CT+RL vs CT-only (if \texttt{--compare-baseline}).
  \item \textbf{Residual torques} $\Delta\tau_1$, $\Delta\tau_2$ [Nm] vs time.
  \item \textbf{Spring extension} $\delta$ [mm] vs time.
\end{enumerate}

Summary statistics (mean, max, std of EE error after the 3\,s preamble)
are printed to stdout.

\subsection{Deterministic Evaluation}

During evaluation, actions are sampled with \texttt{deterministic=True}:
\begin{lstlisting}
action, _ = model.predict(obs, deterministic=True)
\end{lstlisting}
This uses the mean of the Gaussian policy (no exploration noise),
giving the best-case performance of the learned policy.


% =============================================================================
\section{File Structure}
% =============================================================================

\begin{center}
\begin{tabular}{lp{9cm}}
\toprule
\textbf{File} & \textbf{Role} \\
\midrule
\texttt{rl/\_\_init\_\_.py} & Package init \\
\texttt{rl/envs/\_\_init\_\_.py} & Exports \texttt{ManipulatorResidualEnv} \\
\texttt{rl/envs/manipulator\_residual\_env.py} & Gymnasium environment (obs, action, reward, reset, step) \\
\texttt{rl/train\_ppo\_residual.py} & Isaac Sim scene setup + SB3 PPO training loop \\
\texttt{rl/eval\_ppo\_residual.py} & Load checkpoint + run deterministic episode + plot \\
\texttt{rl/checkpoints/} & Saved model weights (\texttt{.zip}) \\
\texttt{rl/logs/} & CSV + TensorBoard training logs \\
\texttt{rl/notes/} & This document \\
\bottomrule
\end{tabular}
\end{center}

Dependencies on project modules:
\begin{itemize}[nosep]
  \item \texttt{robots.cup\_manipulator\_tendon\_isaac} --- robot class, FK, IK
  \item \texttt{controller.computed\_torque\_isaacsim} --- CT controller,
        \texttt{ik\_to\_joint\_space\_references}
  \item \texttt{controller.trajectory} --- \texttt{CircleTrajectory},
        \texttt{RectTrajectory}, \texttt{PreambleTrajectorySource}
  \item \texttt{actuators.sea\_isaacsim} --- \texttt{SEACableActuatorNP},
        \texttt{SEADiagnostics}
\end{itemize}


% =============================================================================
\section{Usage Reference}
% =============================================================================

\begin{lstlisting}[language=bash, caption={Training commands}]
conda activate env_isaacsim

# Default training (headless, circle, 500k steps)
python rl/train_ppo_residual.py

# Soft spring (more SEA lag -- RL has more to learn)
python rl/train_ppo_residual.py \
    --spring-stiffness 100 --motor-bandwidth 20 \
    --total-timesteps 2000000

# Resume from checkpoint
python rl/train_ppo_residual.py \
    --resume rl/checkpoints/ppo_residual_latest

# Monitor training curves
tensorboard --logdir rl/logs
\end{lstlisting}

\begin{lstlisting}[language=bash, caption={Evaluation commands}]
# Evaluate latest model
python rl/eval_ppo_residual.py

# Compare RL vs CT-only baseline
python rl/eval_ppo_residual.py --compare-baseline

# Evaluate with live rendering
python rl/eval_ppo_residual.py --render native
\end{lstlisting}


% =============================================================================
\section{Future Work}
% =============================================================================

\begin{enumerate}
  \item \textbf{Domain randomization} --- randomize $k_s$, $b_c$, $\omega_m$,
        joint friction, and payload mass at each \texttt{reset()} to improve
        transfer robustness.
  \item \textbf{Multi-trajectory training} --- cycle through rect/circle/line
        trajectories across episodes so the policy generalizes beyond a single
        shape.
  \item \textbf{Isaac Lab port} --- move to \texttt{ManagerBasedRLEnv} for
        1024+ parallel environments on GPU, leveraging vectorized PhysX.
  \item \textbf{Observation normalization} --- wrap with SB3's
        \texttt{VecNormalize} to auto-scale observations and rewards.
  \item \textbf{Curriculum learning} --- start with a stiff SEA ($k_s = 500$),
        progressively soften during training.
\end{enumerate}


\end{document}
